\enteteExam{NFP 107 -- Systèmes de gestion de bases de données - Durée 2h30}%
{Seconde session -- 1er septembre 2021}%
{Effectué à distance}%

Nous gérons des collections de photos (numériques) que nous classons dans des albums. 
Notre application est de plus capable d'identifier la présence de personnes sur une
photo. Cela nous donne le schéma relationnel suivant\,:

      \begin{itemize}
        \item Personne (\textbf{idPersonne}, prénom)
        \item Photo (\textbf{idPhoto}, lieu, date, idAuteur, idAlbum)
        \item Album (\textbf{idAlbum}, nom)
        \item  Présence (\textbf{idPrésence}, idPhoto, idPersonne)
       \end{itemize}

Les clés primaires sont en gras, les clés étrangères se déduisent aisément (notez que l'auteur
d'une photo est une personne).


\section{Schéma (8 pts)}

\begin{enumerate}

   \item ({\bf 2 pts}) Donnez les commandes {\tt create table} décrivant
         le schéma des relations ci-dessus. Indiquez soigneusement
         les clés étrangères et les clés primaires.
   \item  ({\bf 2 pts}) Donnez le schéma entité-association correspondant à ce schéma relationnel.
   \item ({\bf 1 pts}) Une des tables résulte-t-elle d'une réification\,? Justifiez.
   \item ({\bf 1 pts}) Est-il possible avec ce schéma de représenter le fait qu'une même personne 
   figure deux fois sur la même photo\,? Si oui, comment pourrait-on améliorer le schéma pour éviter 
   que ce soit possible\,?
 
 \item ({\bf 2 pts}) On voudrait qu'une photo puisse figurer sur plusieurs album. Indiquez quelle
 dépendance fonctionnelle disparaît, et les modifications que cela apporterait au schéma.

\end{enumerate}

\section{Requêtes (8 pts)}

\textbf{Important}\,: ces requêtes sont à exprimer
sur le schéma donné dans l'énoncé, sans modification.

\begin{enumerate}
% 1 pt
  \item Donnez les prénoms des auteurs de photos situées en Corse.
    
% 1 pt
  \item Donnez les prénoms de tous  les auteurs de \emph{selfies} (l'auteur est présent sur la photo).
  
  
% 1 pts
  \item Donnez les noms des albums dont au moins une photo montre la personne prénommée ``Sophie'',
     sans utiliser l'imbrication des requêtes.
  % 1 pts
  \item Même question, mais  en utilisant systématiquement l'imbrication des requêtes.
  
  \item Quelles personnes ne sont auteur d'aucune photo\,?
% 1 pts
  \item Donnez les identifiants des photos de l'album "Bretagne 2021" sur lesquelles ne figure
     pas la personne ``Philippe''.

  \item Donnez le nom des albums et le nombre de photos qu'ils contiennent
  
\item Donnez les noms des albums contenant des photos de plusieurs auteurs différents.
  \end{enumerate}

\section{Algèbre relationnelle ({\bf 2 pts})}

Exprimez les requêtes 3 et 4 précédentes
en algèbre relationnelle.

\section{Modèle relationnel ({\bf 2 pts})}

Soit les attributs P, L, A, N  (P est la photo, L le lieu,
A l'album et N le nom). Les dépendances fonctionnelles sont\,:
\texttt{P -> L} et \texttt{A -> N}

Répondre aux questions suivantes\,: 
\begin{enumerate}
  \item Quelle est la clé\,?
  \item Quel est le résultat de l'algorithme de normalisation\,?
  \item (bonus pour 2 points)\,: ce modèle correspond-il 
     au schéma de l'énoncé\,? Si non quelle est la différence\,?
\end{enumerate}